\section{The Nature, Incorporeality, Place, and Movement of the Angels}\label{sec:st-nature}

How can a creature possess no matter at all---no flesh, no boundary, no weight, no quantitative parts---and yet remain infinitely distant from God, distinct from every other spirit, present in a single room, and able to travel from heaven to Nazareth? St.~Thomas Aquinas opens his systematic treatise on the angels in the \emph{Summa theologiae} (\ST{I, qq.~50--53}) by confronting that exact metaphysical puzzle. Every answer he gives turns on a single principle: an angel is neither a subtle body nor an emanation of the divine substance, but a finite, subsisting intellectual form whose being is received from God.

\subsection{Substance, Immateriality, Multitude, Species, and Incorruptibility (\ST{I, q.~50})}

Aquinas begins with the sheer incorporeality of angelic nature (\ST{I, q.~50, a.~1}). Sacred Scripture declares of God, ``Who makest thy angels spirits: and thy ministers a burning fire'' (Ps~103:4 [Heb~104:4]), and St.~Luke records that the Sadducees erred precisely in denying resurrection, angel, and spirit alike (Acts~23:8). Reasoning from the perfection of the universe, Aquinas argues that an effect achieves its highest perfection when it returns to its principle by resembling the very mode in which its cause produces it. Because God brings creatures into being through intellect and will, the order of creation would remain incomplete without purely intellectual creatures. Yet the act of understanding cannot be the act of any bodily organ, since a body is pinned down to ``here'' and ``now'' while intellect grasps universal truth; therefore there must exist wholly incorporeal intellectual substances (\ST{I, q.~50, a.~1}).

Ancient materialists denied incorporeal spirits because they could not rise above imagination, which pictures only bodies. Even among Christian Fathers, however, phrases occur that seem at first glance to attribute subtle bodies to angels---most notably St.~John of Damascus's remark in \emph{De fide orthodoxa}~(II.3, quoted in Section~\ref{sec:patristic-witness}) that beside the incomparable Deity every creature appears dense and material, and that an angel is an intelligent essence ``in perpetual motion.'' Aquinas resolves these expressions cleanly (\ST{I, q.~50, a.~1, ad~1--3}). Damascene calls angels dense or corporeal only comparatively beside God's absolute simplicity, as lukewarm water feels cold next to fire while containing no coldness in itself. Likewise, calling angels ``in perpetual motion'' signifies uninterrupted actual intellection rather than bodily locomotion, and St.~Ambrose's statement that every creature is circumscribed (\emph{De Spiritu Sancto}~I.7) denotes essential finitude, not containment within physical dimensions.

A sharper metaphysical question follows immediately: even if an angel has no \emph{body}, might it still be composed of \emph{spiritual matter} and form (\ST{I, q.~50, a.~2})? In the thirteenth century this was a living debate among Catholic masters. Following Avicebron's \emph{Fons vitae} and seeking to safeguard the distinction between Creator and creature, St.~Bonaventure defended universal hylomorphism in his commentary on the \emph{Sentences} (\emph{In II Sent.}, d.~3, p.~1, a.~1, q.~1). While ruling out quantitative, heterogeneous, or body-soul composition in angels, Bonaventure extended the name of matter broadly to every constitutive principle of potentiality, concluding \emph{ipsa substantia Angeli composita est ex materia et forma}, ``the substance of the angel itself is composed of matter and form'' (\emph{In II Sent.}, d.~3, p.~1, a.~1, q.~1). Because an angel possesses mutability not merely toward non-being but across properties, passibility, individuation, limitation, and essential composition, Bonaventure saw no other ground to distinguish a subsisting creature from divine simplicity than two constitutive principles related as potential and actual. In his view, Boethius's remark in \emph{De Trinitate} that God alone is immaterial demands spiritual matter in every subsisting creature, since matter considered in itself is neither spiritual nor corporeal until actuated by form.

Aquinas rejects universal hylomorphism on two decisive grounds (\ST{I, q.~50, a.~2}). First, spiritual and corporeal creatures cannot share a single common matter, because spiritual and bodily forms cannot occupy the same part of matter without making one thing both spiritual and bodily; yet matter cannot be divided into distinct parts except through quantity, which would subject spiritual matter to bodily dimensions. Second, every substance acts according to its mode of being. Understanding is wholly immaterial, because forms immersed in matter are individuated and cannot be grasped intellectually until stripped of material conditions. Therefore, in the exact words of the Leonine Latin text (\emph{omnis substantia intellectualis est omnino immaterialis}), ``every intellectual substance is altogether immaterial.'' Pseudo-Dionysius stands on the same ground: ``The first creatures are understood to be as immaterial as they are incorporeal'' (\emph{De divinis nominibus}~IV). When human minds distinguish genus and difference in an angel, they do not uncover matter and form inside the angel; they simply apprehend a higher, simpler substance according to the composite mode of human cognition (\ST{I, q.~50, a.~2}; \emph{De ente et essentia}, ch.~5).

How, then, does an angel without matter avoid collapsing into the divine simplicity? Aquinas answers with the real distinction between essence and existence (\emph{esse}), worked out early in \emph{De ente et essentia} (chs.~4--5) and condensed in \ST{I, q.~50, a.~2, ad~3--4}. A material substance carries a double composition: matter and form constitute its nature, and that composite nature stands as potency to the act of existence by which it actually is. Remove matter, and let the form subsist by itself: the second composition remains intact. Every essence other than God can be conceived without knowing whether it exists in reality---one can understand what a man or a phoenix is while remaining ignorant of whether it exists (\emph{De ente et essentia}, ch.~4). Subsistent existence itself (\emph{esse tantum}) can be only one; in every angel, therefore, the subsisting form---Boethius's \emph{quod est}, ``what is''---stands as potentiality receiving a limited act of existence---Boethius's \emph{quo est}, ``whereby it is''---from the First Cause. Hence the \emph{Liber de causis} (prop.~16) describes an angelic intelligence as ``finite from above,'' because its existence is received and bounded by its nature, and ``infinite from below,'' because its form is uncontracted by any receiving matter (\ST{I, q.~50, a.~2, ad~4}; \emph{De ente et essentia}, ch.~5).

From immateriality Aquinas turns to the multitude of the angels (\ST{I, q.~50, a.~3}). Daniel's vision beholds an innumerable court: ``thousands of thousands ministered to him, and ten thousand times a hundred thousand stood before him'' (Dan~7:10), and Pseudo-Dionysius declares that ``there are many blessed armies of the heavenly intelligences, surpassing the weak and limited reckoning of our material numbers'' (\emph{Celestial Hierarchy}~XIV). Aquinas sets aside three philosophical constrictions: Plato's identification of separate substances with the exemplar species of sensible things; Aristotle's tentative conjecture (\emph{Metaphysics}~XII) restricting separate intelligences to the number of celestial spheres; and Moses Maimonides's attempt to reconcile Aristotle with Scripture by taking scriptural angels as human envoys or natural forces. Because God chiefly intends the perfection of the universe, nobler kinds of creatures exceed lower ones incomparably: just as the incorruptible heavens dwarf corruptible earthly bodies in physical vastness, immaterial substances surpass all material bodies in multitude. That multitude is not quantitative number, which arises from dividing a continuum, but transcendental multitude grounded in formal distinction (\ST{I, q.~50, a.~3, ad~1}); nor does angelic simplicity imply fewness, for nearness to God reduces internal composition rather than the number of subjects (\ST{I, q.~50, a.~3, ad~2}).

Immateriality yields a still more arresting corollary in \ST{I, q.~50, a.~4}: it is impossible for two angels to belong to the same species. Things that agree in species while differing in number---two oaks, two sparrows, two human beings---share a single specific form multiplied in different parcels of designated matter. Without matter, a subsisting form exhausts its entire species, just as whiteness existing apart from white surfaces could only be one whiteness; as Aquinas states in \emph{De ente et essentia} (ch.~4), quoting Avicenna, \emph{quot sunt ibi individua, tot sunt ibi species}, ``as many individuals as there are, so many species.'' Even if spiritual matter were conceded, it could not divide numerically without quantity, and differences of potentiality would produce differences of genus rather than mere numerical distinction. Aristotle observes that within a single species there is no ``first and second'' (\emph{Metaphysics}~III), yet Pseudo-Dionysius teaches that even inside a single angelic order there are first, middle, and last (\emph{Celestial Hierarchy}~X). Angels therefore differ by distinct formal grades of intellectual perfection, as fire formally excels air, rather than by accidental degrees of intensity; and specific multiplication glorifies the Creator more nobly than numerical repetition, which in corruptible creatures exists chiefly to preserve a perishable species (\ST{I, q.~50, a.~4, ad~1--4}).

Among Catholic schools, \ST{I, q.~50, a.~4} remains one of the clearest flashpoints of genuine disagreement. St.~Bonaventure had already rejected the ``as many species as individuals'' thesis before Aquinas wrote the \emph{Summa}, calling it a \emph{praesumtio} lacking support in Scripture or the Fathers, since Scripture assigns many angels to a single office and common operation; Bonaventure instead defends as \emph{sobria et catholica}, ``sober and catholic,'' the view that angels differ personally within one species, multiplying charity's joy in a common society (\emph{In II Sent.}, d.~3, p.~1, a.~2, q.~1). For Bonaventure, individuation springs from the actual conjunction of spiritual matter and form, like many seals pressed into wax (\emph{In II Sent.}, d.~3, p.~1, a.~2, q.~3). On 7~March 1277, Bishop Stephen Tempier's Parisian syllabus censured three propositions touching this point---articles~81, 96, and 191, recorded by Durandus of Saint-Pour\c{c}ain---specifically rejecting the claim that God \emph{cannot} multiply individuals under one species without matter. A generation later, Bl.~John Duns Scotus likewise rejected both Thomist specific uniqueness and Bonaventurian spiritual matter, grounding angelic individuation within a common species in a positive formal difference, \emph{haecceitas} or ``thisness'' (\emph{In II Sent.}, d.~3, q.~6).

Question~50 closes with angelic incorruptibility (\ST{I, q.~50, a.~5}). Natural corruption occurs only when a form is separated from its matter. Because existence belongs to form intrinsically---just as roundness belongs to a circle by itself, whereas a bronze circle can lose its circular shape---a subsisting form without matter carries no inner principle of decay or contrariety and cannot naturally lose its existence. Its very operation confirms its nature, for the intellect grasps timeless, everlasting truth; Pseudo-Dionysius accordingly teaches that intellectual substances ``have unfailing life, being free from all corruption, death, matter, and generation'' (\emph{De divinis nominibus}~IV). When St.~John of Damascus states that angels possess immortality ``by grace, not by nature'' (\emph{De fide orthodoxa}~II.3), Aquinas reads him as speaking of perfect immutability, since ``every change is a kind of death'' (St.~Augustine, \emph{Contra Maximinum}~III) and angels receive immutability of will only in glory (\ST{I, q.~50, a.~5, ad~1}). Similarly, St.~Gregory the Great's teaching that all creatures, angels included, would slip into nothingness unless sustained by God (\emph{Moralia in Iob}~XVI) expresses their radical dependence on the Creator's conserving power, not any intrinsic corruptibility: a necessary creature still has a cause of its necessity (\ST{I, q.~50, a.~5, ad~3}).

\angeldossier{\ST{I, q.~50, aa.~1--5} (Substance of the angels: incorporeality, immateriality, multitude, specific distinction, and natural incorruptibility).}{Church teaching (aa.~1, 5); Common teaching (a.~3); Thomist position (aa.~2, 4).}{Ps~103:4 [Heb~104:4]; Ps~148:2, 4; Dan~7:10; Acts~23:8; Pseudo-Dionysius, \emph{De divinis nominibus}~IV and \emph{De caelesti hierarchia}~X, XIV; St.~Ambrose, \emph{De Spiritu Sancto}~I.7; St.~Augustine, \emph{Contra Maximinum}~III; St.~Gregory the Great, \emph{Moralia in Iob}~XVI; St.~John of Damascus, \emph{De fide orthodoxa}~II.3; Boethius, \emph{De Trinitate} and \emph{De hebdomadibus}; \emph{Liber de causis}, prop.~16; Aristotle, \emph{Physics}~I, IV, VI, \emph{Metaphysics}~I, III, VIII, XI--XIII, \emph{De anima}~I.5; St.~Thomas Aquinas, \emph{De ente et essentia}, chs.~4--5.}{Subtle aerial or fiery bodies entertained by some early Fathers and comparatively by Damascene (\emph{DFO}~II.3), with Augustine leaving the question open in \emph{City of God}~XXI (a.~1); universal hylomorphism in Avicebron and St.~Bonaventure, \emph{In II Sent.}~d.~3, p.~1, a.~1, q.~1 and \emph{Breviloquium}~II.6 (a.~2); Aristotle's restriction of separate substances to the number of celestial movers (\emph{Metaphysics}~XII) and Maimonides's concordism (a.~3); plurality of angels within one species in St.~Bonaventure, \emph{In II Sent.}~d.~3, p.~1, a.~2, q.~1, Paris 1277 arts.~81, 96, 191, and Bl.~John Duns Scotus's \emph{haecceitas} (a.~4).}

\subsection{Angels in Relation to Bodies (\ST{I, q.~51})}

Having established that an angel is not itself a body, Aquinas asks whether angels possess bodies naturally united to them, as human souls do (\ST{I, q.~51, a.~1}). His answer is firmly negative. Because understanding is not a bodily operation, union with a body does not belong to intellectual substance as such; it belongs solely to the lowest intellectual substance, the human soul, whose intellect stands in potentiality and must gather knowledge through bodily senses. Just as the perfection of the universe requires animals endowed with sense, the perfection of the intellectual order requires pure intelligences needing no sensory medium. Explaining contrary patristic phrases, Aquinas observes that Origen (\emph{Peri Archon}~I) was misled by ancient philosophers; St.~Bernard (\emph{Sermones super Cantica}~VI) refers to an assumed instrument for ministry; St.~Gregory's naming an angel a ``rational animal'' (\emph{Homiliae in Evangelia}~X) is metaphorical; and St.~Augustine (\emph{De Genesi ad litteram}~III; \emph{City of God}~VII) reports the Platonists' ``aerial animals'' without adopting the thesis as his own (\ST{I, q.~51, a.~1, ad~1--2}). Nor does animating a body raise a spirit's dignity, for while God gives life effectively as Creator (1~Kgs [1~Sam]~2:6), giving life formally belongs to an incomplete substantial part (\ST{I, q.~51, a.~1, ad~3}).

In his later \emph{Quaestiones disputatae de malo} (q.~16, a.~1), Aquinas revisits the aerial-body question for demons with notable historical candor. Quoting Augustine's remark in \emph{City of God}~XXI that whether demons have bodies naturally united to them makes little difference to Christian faith, Aquinas first rejects the Peripatetic denial of demons (citing demonic speech in unknown tongues) and Plotinus's identification of demons with dead human souls, then dismantles the Platonist aerial body through Aristotelian physics: air is homogeneous, lacks determinate shape or organic parts, and cannot sustain the sense of touch. Yet in his replies (\emph{De malo}~q.~16, a.~1, ad~3, 7), Aquinas freely grants that Pseudo-Dionysius---who often followed Platonic idiom---and Damascene, echoing Origen, may have spoken provisionally as though spirits possessed subtle bodies.

If angels have no natural bodies, what are the visible figures who walk into Abraham's camp, dine under the oak of Mambre, lodge with Lot in Sodom, and escort young Tobias across Media? Aquinas affirms that angels sometimes assume real external bodies (\ST{I, q.~51, a.~2}). Some writers tried to reduce every biblical apparition to an inner imaginative vision, like a prophet's dream. Scripture refutes that reduction: an imaginary vision exists only inside one seer's mind, whereas the three visitors at Mambre were beheld simultaneously by Abraham, Sara, and their servants, just as the two angels at Sodom were seen by Lot and the townsmen, and Raphael was seen by Tobias's entire family (Gen~18--19; Tob~5--12). Only an external body can be seen publicly by many eyes at once. Angels fashion such bodies out of condensed air by divine power so that they take shape and color as clouds do, uniting themselves to the assumed body ``not as its form, nor merely as its mover, but as its mover represented by the assumed movable body'' (\ST{I, q.~51, a.~2, ad~2--3}). Angels assume bodies not for their own benefit, but for ours: to offer a sensible pledge of the intellectual companionship men will share with them in heaven, and, throughout the Old Testament, to foreshadow the supreme visible mission in which the Word of God would truly assume human flesh (\ST{I, q.~51, a.~2, ad~1}).

An assumed body, however, is not a living organism, so an angel cannot exercise strictly vital functions through it (\ST{I, q.~51, a.~3}). Following Aristotle (\emph{De somno et vigilia}~I), operations proper to life---seeing with eyes, digesting food, begetting offspring---belong only to a living subject endowed with a vital soul. Actions that living bodies share with non-living bodies, by contrast, an angel genuinely produces: moving the assumed body locally as a mover moves its instrument, and vibrating air to form articulate speech. No pretence or falsehood enters these apparitions, any more than in Scripture's sensible metaphors, because the bodily limbs signify spiritual powers rather than falsely claiming biological life (\ST{I, q.~51, a.~3, ad~1--4}). When Abraham sets butter, milk, and veal before his guests (Gen~18:8), there is division and consumption of food, not assimilation into living flesh---unlike Christ's genuine post-Resurrection eating in Luke~24. The Archangel Raphael himself lifts the veil for Tobit and his son: ``I seemed indeed to eat and to drink with you but I use an invisible meat and drink, which cannot be seen by men'' (Tob~12:19; \ST{I, q.~51, a.~3, ad~5}). As for the enigmatic ``sons of God'' of Genesis~6:4, Aquinas follows St.~Augustine (\emph{City of God}~XV.23) in identifying them with the pious descendants of Seth marrying the daughters of Cain, ruling out angelic biological paternity while noting, with Augustine's report on incubi and \emph{De Trinitate}~III, that if a demon ever contrived a human birth, it could only transfer human seed rather than generate from its own spiritual substance (\ST{I, q.~51, a.~3, ad~6}).

\angeldossier{\ST{I, q.~51, aa.~1--3} (Angels in relation to bodies: absence of natural bodily union, assumed bodies, and non-vital exercise of bodily actions).}{Church teaching (a.~1); Common teaching (aa.~2--3).}{Gen~6:4; Gen~18--19; Tob~5:7--8, 12:19; 1~Kgs [1~Sam]~2:6; Ps~8:2; Luke~24; Pseudo-Dionysius, \emph{De divinis nominibus}~IV and \emph{De caelesti hierarchia}; St.~Augustine, \emph{De Genesi ad litteram}~III, \emph{De civitate Dei}~VII, XV--XVI, XXI, and \emph{De Trinitate}~III; St.~Gregory the Great, \emph{Homiliae in Evangelia}~X; St.~Bernard of Clairvaux, \emph{Sermones super Cantica}~VI; Aristotle, \emph{De anima}~II, \emph{De somno et vigilia}~I, \emph{Physics}~VIII; St.~Thomas Aquinas, \emph{De malo}~q.~16, a.~1.}{The Platonist hypothesis of aerial bodies naturally united to spirits, entertained by Origen (\emph{Peri Archon}~I) and left open by St.~Augustine (\emph{De civitate Dei}~XXI; cf.~\emph{De malo}~q.~16, a.~1) (a.~1); none recorded on aa.~2--3.}

\subsection{Angels in Relation to Place (\ST{I, q.~52})}

Every night at Compline, the Church prays over her houses: \emph{Angeli tui sancti, habitantes in ea, nos in pace custodiant}, ``Let Thy holy angels who dwell herein, keep us in peace.'' Aquinas opens Question~52 with that liturgical prayer (\ST{I, q.~52, a.~1}). Yet how does an incorporeal spirit dwell inside a house? Aquinas replies that an angel and a stone are said to be ``in a place'' equivocally. A body occupies space by \emph{contact of dimensive quantity}: its three dimensions match the surrounding surface, so the place measures and encloses the body. An angel possesses no dimensive quantity, only \emph{virtual quantity}---the range and intensity of its active power. An angel is therefore in a corporeal place solely by applying its power to that place. Far from being boxed inside four walls, a spirit virtually contains the bodily place it touches, just as the human soul contains and unifies the body rather than being held inside it like water in a jar (\ST{I, q.~52, a.~1}). Boethius's remark in \emph{De hebdomadibus} that incorporeal things are not in place, and Aristotle's analysis of place in \emph{Physics}~IV, govern only dimensive, bodily enclosure.

From this finite power follows the impossibility of simultaneous bilocation (\ST{I, q.~52, a.~2}). God's infinite essence and power reach every creature at once as universal cause, so God is everywhere. Because an angel's nature and power are finite, its immediate operation terminates at one determinate object and therefore at one place, confirming St.~John of Damascus's rule (\emph{De fide orthodoxa}~II.3) that an angel is present and operates mentally at its assigned post without being able to act in several places at once. Aquinas warns here against a seductive trick of the imagination: picturing an angel as an indivisible geometric point pinned to a single point in space. A point is indivisible \emph{within} quantity and position; an angel stands altogether outside quantity and position. Accordingly, an angel may freely apply its power to a tiny body or a vast continuous one, and the entire continuous body immediately touched by its power counts as one place---so even an angel moving a celestial sphere is not everywhere, since its power applies immediately to the sphere's first moved part, and no single angel moves all spheres immediately. This yields the classic threefold scholastic distinction of local presence (\ST{I, q.~52, a.~2}):
\begin{quotation}
For a body is in a place in a circumscribed fashion, since it is measured by the place. An angel, however, is not there in a circumscribed fashion, since he is not measured by the place, but definitively, because he is in a place in such a manner that he is not in another. But God is neither circumscriptively nor definitively there, because He is everywhere. (St.~Thomas Aquinas, \ST{I, q.~52, a.~2})
\end{quotation}

Can two angels occupy the exact same place at the same time? Aquinas answers no (\ST{I, q.~52, a.~3}), and his proof reveals how little his angelology resembles crude spatial mechanics. Two bodies cannot share one place because their dimensions exclude each other; two angels cannot share one place because two complete proximate causes cannot immediately be the total cause of the exact same effect in the same order. Several rowers can pull one boat only because each is a partial mover whose combined strength makes up one single cause. Because an angel is in a place precisely by touching it immediately as a complete virtual container, two angels cannot simultaneously be the immediate container of the same place. By contrast, an angel and a body coexist without hindrance because their modes of presence differ (virtual versus dimensive), and a human soul coexists with a possessing demon in one body because the soul unites to the body as its substantial form, whereas the demon acts upon it extrinsically (\ST{I, q.~52, a.~3, ad~2--3}).

Later scholastic schools vigorously contested this operational account of angelic place. On 7~March 1277, Parisian article~25 censured the unqualified claim that separate substances are in place solely through operation, ``if it be understood that without operation a substance is not in place.'' Taking up that challenge, Bl.~John Duns Scotus argued that an angel does not exist in place necessarily---God could create an angel apart from any corporeal universe---yet possesses a passive capacity for local presence rooted directly in its limited substance, so that an angel can be definitively present in a place without acting upon it, and several angels can coexist in the same definitive place (\emph{In II Sent.}, d.~2, q.~6, nn.~11, 15). Commenting on \ST{I, q.~52, a.~1}, Cardinal Thomas de Vio Cajetan observed that all doctors agree an angel is in place definitively by virtual rather than quantitative contact, yet the exact ground of that presence exceeds the capacity of the human mind (\emph{excedit enim evidentia eius humani ingenii capacitatem}), since Sacred Scripture assures us that angels are somewhere without explaining how spirit touches place.

\angeldossier{\ST{I, q.~52, aa.~1--3} (Angels in relation to place: presence by virtual contact, definitive presence excluding bilocation, and non-co-location of two angels).}{Common teaching (aa.~1--2); Thomist position (a.~3).}{Gen~19:25; Compline Collect (\emph{Visita, quaesumus, Domine}); St.~Augustine, \emph{De Trinitate}~VI; St.~John of Damascus, \emph{De fide orthodoxa}~II.3; Boethius, \emph{De hebdomadibus}; Aristotle, \emph{Physics}~IV, VII--VIII; Cajetan, \emph{In ST}~I, q.~52, a.~1.}{Bl.~John Duns Scotus distinguishes definitive presence from operative contact (\emph{In II Sent.}~d.~2, q.~6; cf.~Paris 1277 art.~25) (a.~1); Scotus and later Franciscan and Suarezian theologians permit two angels to occupy the same place simultaneously by definitive presence (a.~3); none recorded on a.~2.}

\subsection{The Local Movement of the Angels (\ST{I, q.~53})}

If an angel occupies place by applying its power, local movement is simply a succession of virtual contacts with different places (\ST{I, q.~53, a.~1}). Even beatified angels move locally, just as the Creed confesses that Christ's blessed soul descended into hell. Aristotle's demonstration in \emph{Physics}~VI that an indivisible cannot move applies only to quantitative indivisibles moving continuously across divisible magnitude; an angel is neither confined to a point nor restricted to one mode of transit. It can move either \emph{continuously}---applying its power to a divisible space and gradually releasing one part as it engages the next, without its own substance having parts---or \emph{non-continuously}, withdrawing its power altogether from one whole place and applying it immediately to another. Nor does local movement imply imperfection or personal lack in a blessed angel: rather than the motion of a being in potentiality seeking fulfillment, it is the act of an agent already perfected in act, moving solely for our salvation as ``ministering spirits, sent to minister for them who shall receive the inheritance of salvation'' (Heb~1:14; \ST{I, q.~53, a.~1, ad~2--3}).

This twofold mode of motion resolves the ancient scholastic riddle of whether an angel flying from heaven to earth must pass through all the space between (\ST{I, q.~53, a.~2}). Between any two places lie potentially infinite intermediate places. In \emph{continuous} motion, an angel must traverse every intermediate place, because only continuity allows a mover to cross a potentially infinite series without numbering actual infinities. In \emph{non-continuous} motion, however, an angel is not confined by space as a stone is, but stands above space as its virtual container; it lies freely within the angel's choice to quit one place whole and apply itself directly to a distant place without touching the intervening air or sky. In such discontinuous transit, each transition is an indivisible atomic step rather than the endpoint of a preceding continuous slide (\ST{I, q.~53, a.~2, ad~3}).

Yet non-continuous motion does not mean that angelic travel happens in a single instant (\ST{I, q.~53, a.~3}). Several medieval masters argued that angelic locomotion is strictly instantaneous, comparing it to sunlight illuminating the dark air at dawn or fire suddenly informing wood: during the entire preceding span of time the air is dark, and in the final bounding instant it is lit, leaving no last instant of darkness. Aquinas uncovers the subtle flaw in that analogy. Illumination and substantial generation are termini of prior continuous motions---the sun's continuous approach or matter's gradual heating---whereas an angel resting in place~$A$ and then appearing in place~$B$ executes an independent act depending on no preceding continuous motion. Resting throughout a stretch of time means being in place~$A$ at every instant of that stretch, including its final instant; therefore there must be a last ``now'' in which the angel was in~$A$, and a distinct, subsequent ``now'' in which the angel is in~$B$. Wherever distinct ``nows'' succeed one another as before and after, there is time.

Accordingly, an angel's local movement always takes place in time: in continuous time if the movement is continuous, and in \emph{discrete, non-continuous time}---a chain of indivisible angelic instants---if the movement is non-continuous (\ST{I, q.~53, a.~3}). Crucially, neither measure is the solar or celestial clock that governs earthly motion, for angelic operations do not depend on the revolution of the heavens. While angelic substantial being endures in the \emph{aevum}, a creaturely mean between God's timeless eternity and earthly flux, angelic acts of intellect, will, and local transit unfold in their own spiritual succession of instants. In discrete angelic time there is no common measure with continuous physical time, and the swiftness of an angel's arrival is determined not by a mechanical ratio of power to resistance, but solely by the determination of the angel's will (\ST{I, q.~53, a.~3, ad~1, 3}).

\angeldossier{\ST{I, q.~53, aa.~1--3} (Local movement of the angels: possibility of local motion, traversal of intermediate space, and angelic discrete time).}{Common teaching (a.~1); Thomist position (aa.~2--3).}{Apostles' Creed (Christ's descent into hell); Heb~1:14; Aristotle, \emph{Physics}~III--VI.}{Some scholastics require traversal of all intermediate space in angelic locomotion (a.~2); Scotist and Suarezian accounts treat the duration of angelic local motion as continuous rather than discrete angelic time (a.~3); none recorded on a.~1.}
